% ============================================================================
% GEM SECTION 4 --- Discussion.  BOTH PARAGRAPHS ARE THE MAP'S TEXT VERBATIM.
% Nothing is added and nothing is trimmed here.  (An earlier revision of this
% file carried two extra sentences lifted from the full-paper discussion --
% the "not unique to COMPOSE" concession and the exact-backward-values
% qualification.  Those were not in the map and have been removed.  If they
% are wanted back, that is an authoring decision, not a fitting decision.)
% ============================================================================
\section{Discussion}
\label{sec:discussion}

COMPOSE is built around a reference process that is learned independently of any
downstream objective. In our experiments, the same GuacaMol-trained $R_\theta$ is
reused for variable-size editing on ZINC-250k and for protein-specific lead
optimization, with only the controller changing between tasks. Representing the
intermediate states explicitly as molecules also makes it possible to look ahead
before committing an edit, continue from a molecule already reached when the
objective changes, and restrict which molecular states can be visited along a
trajectory.

The current formulation still has important limits. Exact finite-horizon control
requires exact backward values, which are only tractable on small enumerable
state spaces. At molecular scale we instead rely on the learned approximation
$h_\phi$, so the exact terminal-law guarantee no longer applies. COMPOSE is also
limited by the rewrite vocabulary used to define its support. Chemistry that
cannot be expressed by those rewrites cannot be recovered through control.
Finally, the optimization experiments inherit the limitations of the property
models and docking scores used as objectives. Extending the rewrite process to
broader chemistry and testing it against experimentally measured objectives are
natural next steps.
